\subsection{A thirteen-fold cover of a second characteristic-two curve}

\begin{claim}{A.C2.4}{Proved}
Let $K$ be finitely generated over $\F_2$. The geometrically
prime-to-two point-section map is injective for the curve
\[
                         C_0:\quad y^2+y=x^3+\frac1x,         \tag{C2.12}
\]
and for its finite-extension forms and hyperbolic open subcurves.
There is a connected cyclic \'etale cover of degree thirteen with
two jointly point-injective quotient maps to a genus-six curve of
$2$-rank zero.
\end{claim}
\begin{proof}
The change $u=1/x$, $V=u^2y$ gives
\[
                    V^2+u^2V=u^5+u.                         \tag{C2.13}
\]
Write $P_\infty$ for the unique point above $u=\infty$ and
$P_0$ for that above $u=0$. In this model the Artin--Schreier
pole orders are one at infinity and three at zero. In particular,
$C_0$ has genus two and $2$-rank one. One can see the latter
from its differential basis $du/u^2,du/u$: Cartier kills the
first and fixes the second.

We first exhibit a thirteen-torsion point on a specified translate
of the curve in its Jacobian. Choose
\[
 \alpha^2+\alpha+1=0,\qquad
 \beta^2+\alpha^2\beta=1,
\]
so that $\alpha,\beta\in\F_{16}$, and set $R=(\alpha,\beta)$
on (C2.13). Define
\[
 \begin{split}
 B(u)&=u^4+\alpha u^3+u+1,\\
 A(u)&=\alpha(1+\beta)u^6+\alpha^2\beta u^5+\alpha u^4
       +\alpha(1+\beta)u^3\\
     &\hspace{12mm}+(1+\alpha\beta)u^2+\alpha u+\alpha^2.
 \end{split}
\]
Direct expansion in this field gives
\[
 A^2+u^2AB+(u^5+u)B^2=(u+\alpha)^{13}.                        \tag{C2.14}
\]
For $G=A+BV$, one has $G(R)=0$ and
$G(\iota R)=\alpha\ne0$. Its pole order at $P_\infty$ is
thirteen: $A$ has pole order twelve, while $BV$ has pole
order thirteen. Taking the norm in (C2.14), with the conjugate
factor a unit at $R$, therefore gives
\[
                    \operatorname{div}(G)=13R-13P_\infty.
\]
The class $L=\mathcal O(R-P_\infty)$ is nontrivial, since a
function with one simple pole would make the curve rational.
It has exact order thirteen. The cover $D$ defined by $z^{13}=G$
is \'etale and connected; otherwise $G$ would be a thirteenth
power, again giving a function with divisor $R-P_\infty$.
Its genus is fourteen.

Here is the complete Frobenius calculation which controls the
$2$-rank of this cover. On a smooth curve in characteristic two,
let $\mathcal B$ be the image of
$d:F_*\mathcal O_C\to F_*\Omega_C^1$, with the module structure
induced by absolute Frobenius. The sequence
\[
 0\longrightarrow\mathcal O_C\longrightarrow F_*\mathcal O_C
   \longrightarrow\mathcal B\longrightarrow0                \tag{C2.15}
\]
is exact. Locally in a parameter $t$, the image is the free
rank-one module of differentials $h(t)^2dt$, so $\mathcal B$
is a line bundle. A rational locally exact differential with
ordinary divisor $2D_1$ is a rational section of $\mathcal B$
with divisor $D_1$; this follows from the same square-coefficient
local description.

For (C2.13), the different formula gives
\[
 \operatorname{div}(du)=4P_0-2P_\infty,
 \qquad \operatorname{div}(u)=2P_0-2P_\infty.
\]
The locally exact differential $d(1/u)=du/u^2$ therefore has
divisor $2P_\infty$. Consequently
\[
                         \mathcal B\simeq\mathcal O(P_\infty).
\]
Tensor (C2.15) by $L$ and use the projection formula. The map
on first cohomology is the semilinear Frobenius map
\[
                   H^1(C,L)\longrightarrow H^1(C,L^2).
\]
Both spaces have dimension one, by Riemann--Roch and nontriviality
of these degree-zero line bundles. Its kernel contains, and equals
the image of, $H^0(C,\mathcal B\otimes L)=H^0(C,\mathcal O(R))$,
which has dimension one. Indeed the preceding $H^0(C,L^2)$
is zero. Thus this Frobenius map is zero.

Over an algebraic closure the eigenspace decomposition for the
cyclic cover is
\[
 H^1(D,\mathcal O_D)=\bigoplus_{i\in\Z/13}H^1(C,L^i).
\]
Each nonzero-index summand is one-dimensional. Frobenius sends
the summand indexed by $i$ to that indexed by $2i$. The powers
\[
                1,2,4,8,3,6,12,11,9,5,10,7
\]
are all the nonzero residues modulo thirteen. Thus every
nonzero-index summand lies on a single Frobenius cycle, containing
the zero map just found. The twelfth Frobenius power is zero on
the sum of these summands. The trivial summand is
$H^1(C,\mathcal O_C)$ and has stable Frobenius rank one.
The Artin--Schreier cohomology sequence identifies $2$-rank with
the dimension of Frobenius-fixed cohomology
\citep[Definition 1.1 and the following discussion]{A_Yang2016},
so $D$ has $2$-rank one as well.

The norm identity lifts the hyperelliptic involution to
\[
                  \tau:z\longmapsto (u+\alpha)/z.
\]
If $\delta$ generates the cyclic deck group, then
$\tau\delta\tau=\delta^{-1}$. The quotient $B=D/\langle\tau\rangle$
has genus six. Indeed, above each of the two branch points of
$C\to\mathbf P^1$, exactly one of the thirteen points is fixed
by $\tau$. Their different exponents are respectively four and
two, unchanged by the \'etale local base change
$D\to C$. Riemann--Hurwitz gives
$26=2(2g(B)-2)+6$, hence $g(B)=6$.

The same dihedral representation argument as in A.C2.2 gives
\[
                           J_D\sim J_C\times J_B^2.
\]
For clarity, the complement of the cyclic invariants in a rational
Tate module has dimension twenty-four. Its fixed subspace under
$\tau$, and its translate by $\delta$, each have dimension twelve.
Their intersection is zero, since an element in it is fixed by
$\delta^2$ and therefore by $\delta$, whose order is odd.
These are precisely the two pulled-back copies of $VJ_B$.
The two maps from $J_B$ thus give the claimed isogeny on the Prym.
Additivity of $2$-rank and the calculation $f(D)=f(C)=1$ imply
$f(B)=0$.

Two deck-conjugate quotient maps $D\to B$ jointly separate points:
a nontrivial collision would again give a fixed point of
$\delta^2$, as in A.C2.3. After composing with an Abel embedding
into $J_B$, their targets have no nonzero geometric $2$-primary
torsion points. The section-lifting argument of A.10 proves the
claimed injectivity. All equations, the deck transformation and
chosen base points are defined after finite extensions; this proves
the stated field and form variants.
\end{proof}
